\section{Integridad}

\begin{caso}{Intrusión y modificación no autorizada de productos de AVAST (CCleaner)}
Leer el caso de intrusión y modificación no autorizada de productos de AVAST.
\end{caso}

\begin{fuentes}
  \fuente{https://www.forbes.com/sites/thomasbrewster/2019/10/21/avast-hacked-again-as-spies-steal-its-passwords/}
  \fuente{https://thehackernews.com/2018/04/ccleaner-malware-attack.html}
  \item \footnotesize Otros disponibles en la web.
\end{fuentes}
\nocite{avast-forbes,ccleaner-thn}

\pregunta{¿Cuál era el objetivo final del atacante?}
\begin{respuesta}
\end{respuesta}

\pregunta{¿Cómo se logró la modificación de los datos o servicios?}
\begin{respuesta}
\end{respuesta}

\pregunta{¿Qué mecanismo(s) pudo haberse implementado para evitar la intrusión y la
modificación no autorizada?}
\begin{respuesta}
\end{respuesta}

\subsection*{Validación de entendimiento del caso}

\begin{mcq}{En el caso de CCleaner estudiado previamente, ¿cuál fue el objetivo más
probable del atacante? (1 o más respuestas válidas):}
  \opcion{Implantar código malicioso dentro del código fuente del aplicativo CCleaner
    con el objetivo de que cada persona que descargue CCleaner también descargue el
    código malicioso y se infecte.}
  \opcion{Eliminar el código fuente de CCleaner para causar un impacto operativo a Avast.}
  \opcion{Ingresar de manera no autorizada a la red de Avast para copiar secretos
    comerciales.}
  \opcion{Comprobar que una empresa de seguridad no está exenta de tener
    vulnerabilidades.}
\end{mcq}

\porque{verificar contra las diapositivas}{%
  Igual que en la sección anterior, las opciones \emph{b}, \emph{c} y \emph{d} llegaron
  fusionadas por el OCR (\texttt{ws-01.md:44-46}); el corte se reconstruyó por sentido.
  Coteja contra la diapositiva original antes de entregar.%
}
